\documentclass[11pt, a4paper]{ctexart}

\usepackage{amsmath, amssymb, amsthm}
\usepackage{geometry}
\usepackage{fancyhdr}
\usepackage{enumitem}
\usepackage{xcolor}
\usepackage{hyperref}
\usepackage{fontspec}
\setCJKmainfont{Songti SC}
\setCJKsansfont{PingFang SC}
\setCJKmonofont{Songti SC}

\geometry{margin=2.2cm}
\pagestyle{fancy}
\fancyhf{}
\rhead{Math 1A}
\lhead{Doubao}
\rfoot{Page \thepage}

\newcommand{\problem}[1]{\subsection*{Problem #1}}
\newcommand{\solution}{\par\noindent\textbf{Solution.}\par\smallskip}
\definecolor{anscolor}{RGB}{0, 80, 160}
\definecolor{badcolor}{RGB}{180, 0, 0}

\begin{document}

\begin{center}
  {\LARGE Worksheet 4}\\[0.4em]
  {\large Math 1A}\\[0.25em]
  Doubao \quad | \quad 2026-10-09
\end{center}

\hrule
\vspace{0.8em}
\problem{Q1}

Suppose $f$ is a function and $L$ and $a$ are real numbers.

\textit{\color{badcolor}\footnotesize 未自动求解: 概念题需要 LLM 推理（可接入 Qwen/Kimi API）}

\noindent\textbf{(a)}
\textit{\color{badcolor}\footnotesize 概念题需要 LLM 推理（可接入 Qwen/Kimi API）}

\noindent\textbf{(b)}
\textit{\color{badcolor}\footnotesize 概念题需要 LLM 推理（可接入 Qwen/Kimi API）}

\noindent\textbf{(c)}
\textit{\color{badcolor}\footnotesize 概念题需要 LLM 推理（可接入 Qwen/Kimi API）}

\noindent\textbf{(d)}
\textit{\color{badcolor}\footnotesize 概念题需要 LLM 推理（可接入 Qwen/Kimi API）}

\vspace{0.6em}\hrule\vspace{0.6em}
\problem{Q2}

Is $\infty$ a number? What does $\lim_{x \to a} f(x) = \infty$ mean? Does $\lim_{x \to a} f(x) = \infty$ mean that the limit exists?

\solution
$\infty$ **不是实数**。

$\lim_{x\to a} f(x) = \infty$ 表示 $f(x)$ 在 $x\to a$ 时无界增长（发散），这个极限**在通常意义下不存在**。

\noindent{\bfseries Final Answer:}
\[\boxed{\textcolor{anscolor}{\infty 不是数；该极限为发散}}\]

\vspace{0.6em}\hrule\vspace{0.6em}
\problem{Q3}

What does the Squeeze Theorem say? Draw a picture illustrating this theorem.

\solution
**夹逼定理 (Squeeze Theorem):** 若 $g(x) \le f(x) \le h(x)$ 在 $a$ 附近成立，且 $\lim_{x\to a} g(x) = \lim_{x\to a} h(x) = L$，则 $\lim_{x\to a} f(x) = L$。

几何直观：$f$ 被上下两个函数"夹住"，两者都收敛到 $L$，$f$ 也必收敛到 $L$。

\noindent{\bfseries Final Answer:}
\[\boxed{\textcolor{anscolor}{L（由夹逼定理）}}\]

\vspace{0.6em}\hrule\vspace{0.6em}
\problem{Q4}

What do we mean by $\lim_{x \to a^-} f(x) = L$? What is meant by $\lim_{x \to a^+} f(x) = L$? Graph a function $y = f(x)$ and some point $x = a$ where $\lim_{x \to a^-} f(x) \neq \lim_{x \to a^+} f(x)$. Find an algebraic expression for such a function.

\solution
$\lim_{x \to a} f^{2} x^{2} = a^{2} f^{2}$

\noindent{\bfseries Final Answer:}
\[\boxed{\textcolor{anscolor}{a^{2} f^{2}}}\]

\vspace{0.6em}\hrule\vspace{0.6em}
\problem{P1}

(a) Is the limit of a sum always the sum of the limits? Give an example where $\lim_{x \to a}[f(x) + g(x)]$ exists even though neither $\lim_{x \to a} f(x)$ nor $\lim_{x \to a} g(x)$ exist.

\solution
（无法直接计算：'list' object has no attribute 'has'）

$\lim_{x \to a} f x = a f$

$\lim_{x \to a} g x = a g$

\noindent{\bfseries Final Answer:}
\[\boxed{\textcolor{anscolor}{a g}}\]

\vspace{0.6em}\hrule\vspace{0.6em}
\problem{P2}

(a) Find $\lim_{x \to 0^+}(\frac{1}{x} - \frac{1}{|x|})$, $\lim_{x \to 0^-}(\frac{1}{x} - \frac{1}{|x|})$, and $\lim_{x \to 0}(\frac{1}{x} - \frac{1}{|x|})$, if they exist.

\solution
$\lim_{x \to 0} - \frac{1}{\left|{x}\right|} + \frac{1}{x} = 0$

$\lim_{x \to 0} - \frac{1}{\left|{x}\right|} + \frac{1}{x} = -\infty$

$\lim_{x \to 0} - \frac{1}{\left|{x}\right|} + \frac{1}{x} = 0$

\noindent{\bfseries Final Answer:}
\[\boxed{\textcolor{anscolor}{0}}\]

\vspace{0.6em}\hrule\vspace{0.6em}
\problem{P3}

For what kinds of functions can one evaluate the limit $\lim_{x \to a} f(x)$ by just plugging in $x = a$? Give an example of a function $f(x)$ for which the limit $\lim_{x \to a} f(x)$ exists and can be evaluated, yet for which this substitution doesn't work.

\textit{\color{badcolor}\footnotesize 未自动求解: 物理题: 无法识别具体公式，请补充已知量}

\vspace{0.6em}\hrule\vspace{0.6em}
\problem{AP1}

Below are listed three limits. Match the best method to each limit and evaluate each limit.

\solution
\noindent{\bfseries Final Answer:}
\[\boxed{\textcolor{anscolor}{\text{见子题}}}\]

\noindent\textbf{(a)}
$\lim_{x \to 0} \frac{x^{1000} - 703 x^{506} + \pi x^{17} + 480}{- 372 x^{75} + 39 x^{14} + \sqrt{7} x^{5} - 24} = -20$
$\boxed{-20}$

\noindent\textbf{(b)}
$\lim_{x \to 4} \frac{x^{2} + 2 x - 24}{x - 4} = 10$
$\boxed{10}$

\noindent\textbf{(c)}
$\lim_{x \to 0} x^{2} \sin{\left(\frac{1}{x} \right)} = 0$
$\boxed{0}$

\vspace{0.6em}\hrule\vspace{0.6em}
\problem{AP2}

Using a calculator, try to guess $\lim_{n \to \infty} n \sin(\frac{\pi}{n})$. (Remember that $n$ is in radians, not degrees!)

\solution
$\lim_{n \to \infty} n \sin{\left(\frac{\pi}{n} \right)} = \pi$

\noindent{\bfseries Final Answer:}
\[\boxed{\textcolor{anscolor}{\pi}}\]

\vspace{0.6em}\hrule\vspace{0.6em}
\problem{AP3}

Draw a circle of radius 1. Inscribe a regular $n$-gon inside it.

\textit{\color{badcolor}\footnotesize 未自动求解: 概念题需要 LLM 推理（可接入 Qwen/Kimi API）}

\noindent\textbf{(a)}
\textit{\color{badcolor}\footnotesize 概念题需要 LLM 推理（可接入 Qwen/Kimi API）}

\noindent\textbf{(b)}
\textit{\color{badcolor}\footnotesize 概念题需要 LLM 推理（可接入 Qwen/Kimi API）}

\noindent\textbf{(c)}
\textit{\color{badcolor}\footnotesize 概念题需要 LLM 推理（可接入 Qwen/Kimi API）}

\vspace{0.6em}\hrule\vspace{0.6em}
\vspace{1em}\noindent\textit{Auto-generated by 国产大模型作业流水线: 7/10 solved.}

\end{document}
